% !TeX spellcheck = it_IT

\documentclass[format=169, 10pt, handout]{beamer}

\usepackage[utf8]{inputenc}
\usepackage[italian]{babel}
\usepackage{hyperref}
\usetheme{metropolis}

\usepackage{mathtools}
\DeclarePairedDelimiter{\ceil}{\lceil}{\rceil}

\title{Esercitazioni di Informatica B}
\subtitle{Funzioni in MATLAB}

\author{Stefano Cereda\\
	stefano.cereda@polimi.it
}
\date{3/12/2018}
\institute[PoliMi]{\vspace{0.5cm}\centering Politecnico di Milano \\ \vspace{0.2cm}
	\includegraphics[width=0.2\textwidth]{../logopolimi}}

%\setbeamercovered{invisible}

\usepackage{xcolor}
\usepackage{listings}

% for matlab code
\definecolor{mygreen}{RGB}{28,172,0} % color values Red, Green, Blue
\definecolor{mylilas}{RGB}{170,55,241}
\definecolor{mGray}{rgb}{0.5,0.5,0.5}
\definecolor{mPurple}{rgb}{0.58,0,0.82}
\definecolor{backgroundColour}{rgb}{0.95,0.95,0.92}

\lstdefinestyle{matlab}{language=Matlab,%
	%basicstyle=\color{red},
	breaklines=true,%
	morekeywords={matlab2tikz},
	keywordstyle=\color{blue},%
	morekeywords=[2]{1}, keywordstyle=[2]{\color{black}},
	identifierstyle=\color{black},%
	stringstyle=\color{mylilas},
	commentstyle=\color{mygreen},%
	showstringspaces=false,%without this there will be a symbol in the places where there is a space
	numbers=left,%
	numberstyle={\tiny \color{black}},% size of the numbers
	numbersep=9pt, % this defines how far the numbers are from the text
	emph=[1]{for,end,break},emphstyle=[1]\color{red}, %some words to emphasise
	%emph=[2]{word1,word2}, emphstyle=[2]{style},    
}

\makeindex


\begin{document}
	\begin{frame}
	\maketitle
\end{frame}

\begin{frame}[fragile]{Funzioni}
Definizione di funzione:
\begin{lstlisting}[style=matlab]
function [r]=pow(x,y)
  r=1;
  for ii=[1:y]
    r = r*x;
  end
end
\end{lstlisting}
\begin{itemize}
	\item lista dei \emph{risultati} fra parentesi quadre
	\item lista degli \emph{argomenti} fra parentesi tonde
	\item realizza una singola operazione che può essere usata in altri programmi
\end{itemize}
\end{frame}

\begin{frame}{Funzione split}
Scrivere uno script Matlab che, ricevuto in ingresso un vettore di interi $m$, stampi due vettori $p$ ed $n$ contenenti i numeri di $m$ positivi e negativi (rispettivamente.)

\pause
Scrivere uno script Matlab che, ricevuto in ingresso un vettore di interi $m$ ed un numero $x$, stampi due vettori $v_1$ e $v_2$ contenenti i numeri di $m$ maggiori e minori di $x$ (rispettivamente.)

\pause
Scrivere una funzione (split) che data una variabile di ingresso $m$ (scalare, vettore o
matrice di numeri) ed un numero $x$ restituisce due vettori $vMin$ e $vMax$. $vMin$
contiene tutti gli elementi di $m$ minori di $x$, mentre $vMax$ gli elementi maggiori
di $x$.
\end{frame}

\begin{frame}{Conversioni}
Scrivere uno script Matlab che, ricevuto un numero $a$, lo converta in binario utilizzando l'algoritmo delle divisioni ripetute.

Modificare lo script per ottenere una funzione che converta in una base a piacere.
\end{frame}

\begin{frame}{Istogramma}
Scrivere una funzione Matlab che, ricevuto in ingresso un array numerico $values$ e due interi $min$ e $max$, restituisca il numero di elementi di $values$ compresi fra $min$ e $max$.

\pause
Scrivere una funzione Matlab che, ricevuto in ingresso un array numerico $values$, ed un intero $bins$ restituisca un array che rappresenti l'istogramma dei valori di bins.
In particolare, l'intervallo di valori contenuti in $values$ andrà diviso in $bins$ intervalli ed ogni intervallo avrà un elemento nel nuovo array contenente il numero di elementi di $values$ che cadono nell'intervallo.

\pause
Se avessimo utilizzato degli script avremmo condiviso il workspace fra i due pezzi di codice, questo potrebbe portare a cancellare delle variabili.
Con le funzioni dobbiamo avere un approccio più pulito in cui ogni funzione lavora solo con i parametri d'ingresso e ritorna tutti i valori di cui potremmo avere bisogno.
\end{frame}

\iffalse
\begin{frame}[allowframebreaks]{Struct e funzioni: codice IBAN}
Un codice IBAN è una sequenza di 27 caratteri alfanumerici così composta:
\begin{itemize}
\item 2 caratteri maiuscoli rappresentanti la sigla internazionale della Nazione
\item 2 cifre (CIN Europeo)
\item 1 carattere maiuscolo (CIN Italiano)
\item 5 cifre (ABI)
\item 5 cifre (CAB)
\item 12 cifre (Numero di Conto)
\end{itemize}
\begin{tabular}{|c|c|c|c|c|c|}
	\hline
	naz & cin\_e & cin\_I & abi & cab	&	conto \\
	\hline
	IT&02&L&12345&12345&123456789012\\
	\hline
\end{tabular}

\framebreak
Si scrivano le seguenti funzioni:
\begin{itemize}
\item \emph{removeSpaces}, che prenda in ingresso una stringa strIn e restituisca una stringa strOut contenente tutti i caratteri di strIn ad esclusione degli spazi.
\item \emph{isAlpha}, che prenda in ingresso una stringa str e restituisca 1 quando str contiene esclusivamente lettere maiuscole, o 0 altrimenti.
\item \emph{isDigit}, che prenda in ingresso una stringa str e restituisca 1 quando str contiene esclusivamente caratteri corrispondenti a cifre, o 0 altrimenti.
\end{itemize}

Si utilizzino tali funzioni per scrivere un’ulteriore funzione \emph{checkIBAN} che riceve come parametro una stringa e restituisce 1 se la stringa inserita, una volta rimossi gli spazi, corrisponde ad un codice IBAN che rispetta lo schema sopra illustrato, o 0 altrimenti.
\end{frame}
\fi

	
\begin{frame}[allowframebreaks]{Struct e funzioni: biblioteca}
Si vuole automatizzare la gestione di una biblioteca di un piccolo paese.

Come primo passo si vuole gestire una matrice di dimensioni $N \times 2$ contenente nella
prima colonna un identificativo numerico e nella seconda colonna un numero di
tessera.

Immaginate l’identificativo numerico come un codice interno alla biblioteca mentre il
numero di tessera come il numero di una tessera esistente appartenente (ad esempio
il vostro codice persona del politecnico).

\begin{tabular}{|c|c|}
	\hline
	ID	&	Codice Persona \\
	\hline
	1	&	123456789	\\
	\hline
	2	&	987654321	\\
	\hline
	...	&	... \\
	\hline
\end{tabular} 


Si forniscano le funzioni necessarie all'inserimento di una nuova tessera e alla ricerca
dell'identificativo numerico dato un numero di tessera (l’identificativo numerico è un
numero crescente che parte da 1).

Funzione \emph{inserisci}:\\
\emph{input}: matrice Soci e nuovo numero di tessera\\
\emph{output}: la matrice Soci

Funzione \emph{cerca}:\\
\emph{input}: matrice Soci e numero di tessera da cercare\\
\emph{output}: l’identificativo numerico associato al numero tessera
\end{frame}

\begin{frame}[allowframebreaks]{Badge - TdE 2 Febbraio 2017}
Un'azienda ha installato un sistema che registra gli ingressi e le uscite dei propri dipendenti tramite
l'uso di badge.
All'entrata e all'uscita, ogni dipendente deve registrarsi passando il badge nel lettore
posto all'ingresso.
Siete stati incaricati di scrivere in Matlab parte del software per la gestione dei
dati acquisiti.

\framebreak

Il programma Matlab è costituito da una serie di funzioni che lavorano sui seguenti dati:
\begin{enumerate}
\item Informazioni relative a ciascun dipendente. Queste sono organizzate in una struct con i
seguenti campi:
\begin{itemize}
\item ID (un numero intero che identifica univocamente ogni dipendente),
\item nome (stringa)
\item  cognome (stringa)
\end{itemize}
\item Il registro dei dipendenti. Questo è un array di struct che contiene in ogni suo elemento:
\begin{itemize}
\item I campi di un dipendente (ID, nome, cognome).
\item La data che corrisponde a un ingresso e a un’uscita del dipendente (si suppone che
avvengano nello stesso giorno).
 Questa data è rappresentata con un numero intero di 8 cifre. Per esempio, il numero 20171223 corrisponde alla data 23 Dicembre 2017.
\item oraIngresso e oraUscita, che vengono rappresentati con un numero intero che indica
l’ora come il numero di minuti trascorsi dalla mezzanotte precedente (ad esempio, le
ore 7:30 corrispondono a 450 minuti, infatti 7 * 60 + 30 = 450).
\end{itemize}
\end{enumerate}

\framebreak

\begin{enumerate}
	\item Per verificare se il dipendente può uscire oppure no, scrivere la funzione \emph{verificaUscita} che prende in ingresso il registro, una data nel formato indicato al punto 2, e un ID e controlla se il
dipendente corrispondente ha lavorato nella data passata come parametro per almeno 7 ore e 30
minuti (i.e., 450 minuti). La funzione restituisce vero se questa condizione risulta verificata, falso in
caso contrario.

\item Scrivere la funzione \emph{badge} che prende in ingresso un dipendente, la struttura registro, la data
odierna (rappresentata come indicato al punto 2), un intero contenente l’ora corrente (rappresentata
come indicato al punto 2).
La funzione \emph{badge} aggiorna registro nel seguente modo: se nella data odierna il dipendente non
compare nella struttura registro (ossia l’utente sta entrando in azienda), la funzione ne aggiunge le
informazioni indicando come data e oraIngresso la data e l’ora passate come parametro.
Altrimenti, se il dipendente è già presente nel registro alla data odierna (ossia l’utente sta uscendo
dall'azienda), la funzione badge aggiorna il campo oraUscita relativo al dipendente con il valore
del parametro ora. La funzione \emph{badge} quindi richiama la funzione \emph{verificaUscita} descritta al punto 1 e, se il dipendente può uscire stampa a video `Arrivederci!', altrimenti `Errore uscita!'

\item  Si supponga che il registro sia stato precedentemente popolato con i dati relativi ai dipendenti di
un'azienda. Si scriva un frammento di codice per definire una variabile \emph{dip1} corrispondente al
dipendente “Giacomo Puccini” avente ID = 1858. Si invochi inoltre la funzione badge al fine di
inserire nel registro il suo ingresso il 3 Febbraio 2017 alle 8:15.
\end{enumerate} 
\end{frame}

\begin{frame}[allowframebreaks]{Atleti - TdE 16 Febbraio 2018}
Un comitato sportivo di sci vuole memorizzare le performance degli sportivi che percorrono la più ripida  
delle piste del comprensorio.
Per ogni atleta, la memorizzazione parte da quota $h_0 = 3000 m$ e termina a  
quota $h_{end} = 1000 m$, e viene effettuata in 100 punti diversi del tracciato equidistanti tra loro.
I dati vengono salvati in:
\begin{itemize}
	\item la matrice $atleti$, le cui righe rappresentano i diversi atleti monitorati, e le cui colonne rappresentano i vari tempi collezionati sul tracciato (quindi l’elemento $(i,j)$ della matrice contiene il tempo in cui l'atleta $i$-­esimo è transitato per il$j$-­esimo punto della pista);
	\item il vettore colonna $velocitaMedia$, con stesso numero di righe della matrice $atleti$, che contiene un valore numerico per ogni atleta, indicativo della sua velocità media lungo tutta la pista.
\end{itemize}  

\framebreak

Scrivere in linguaggio Matlab una funzione $dividiMatrice$ che:
\begin{enumerate}
	\item riceva in input una matrice $atletiTot$ (con la stessa struttura di $atleti$), un vettore $velocitaM$ (con stessa struttura di vettore $velocitaMedia$) e uno scalare soglia;
	\item fornisca in output due matrici $atletiOver$ e $atletiUnder$ (ognuna con la stessa struttura di $atletiTot$).
	$atletiOver$ include solo le righe di $atletiTot$ corrispondenti agli elementi di $velocitaM$ con valore maggiore o uguale di soglia.
	$atletiUnder$, invece, include solo le righe di $atletiTot$ corrispondenti agli elementi di $velocitaM$ con valore minore di soglia. 
\end{enumerate}
\end{frame}

\begin{frame}{Matrice simmetrica}
Scrivere una funzione per controllare se una matrice è simmetrica.	
\end{frame}

\begin{frame}[allowframebreaks]{Partizione - TdE 4 Febbraio 2016}
Si implementi in linguaggio MATLAB una funzione chiamata \emph{partition} che prende in ingresso un  
vettore $V$ di valori numerici non ordinati, e restituisce in uscita un valore numerico $m$ (valore  
centrale di $V$) e un nuovo vettore $T$, definiti nel modo seguente:  

\framebreak

\begin{itemize}
	\item Se $V$ ha lunghezza dispari, $m$ viene posto pari all'elemento in posizione centrale di $V$, e il  
vettore $T$ viene costruito dagli elementi di $V$, disponendo a sinistra di $m$ tutti gli elementi minori  
dello stesso $m$ e a destra tutti gli elementi maggiori di $m$.  
Ad esempio, dato il vettore $V = [6 5 3 4 3 1 2]$, la funzione restituisce il valore centrale $m = 4$ e  
il vettore $T = [ 3 3 1 2 4 6 5]$.  

	\item Se $V$ ha lunghezza pari, il valore $m$ è definito come la media dei due elementi centrali di $V$, e gli  
elementi di $V$ vengono disposti in $T$ in modo tale che tutti quelli minori di $m$ precedono tutti  
quelli maggiori dello stesso numero. In questo caso, non essendo $m$ parte di $V$, esso non viene  
inserito in $T$. 
Ad esempio, dato il vettore $V = [5 6 1 2 3 4]$, la funzione restituisce il valore centrale $m = 1.5$  
(cioè la media dei due valori in posizione centrale 1 e 2) e il vettore $T = [1 5 6 2 3 4]$. 
\end{itemize}
\end{frame}


\end{document}
